% =====================================================================
% Methods subsection: school clustering methodology (Week-5 / v3)
% Paste into the Methods section of the main paper (Overleaf), after the
% rigor-tiering subsection (methods_rigor_tiering.tex).
%
% CITATION KEYS:
%   NEW  (BibTeX provided at the bottom of this file -- paste into your .bib):
%        macqueen1967, ward1963, tibshirani2001, rousseeuw1987,
%        hubert1985, jolliffe2016
%   EXISTING (already cited in the paper -- rename to YOUR keys):
%        nardo2008 -> the OECD composite-indicator handbook (also used in
%                     methods_rigor_tiering.tex)
%
% All numbers below are from the v3 run (clustering_v3_2026-07-24.csv,
% docs/CLUSTERING.md). If the pipeline is re-run, regenerate them.
% Uses natbib (\citep/\citet). Verify years/pages before final submission.
% =====================================================================

\subsection{Grouping Similar Institutions: Feature Engineering and Clustering}\label{sec:clustering}

Beyond the ordinal rigor tiers, the platform calls for an unsupervised grouping of similar
institutions across location, academic profile, and financial characteristics. We deliberately
construct this grouping as a check on the rigor index rather than as a derivative of it: the
rigor score itself is excluded from the feature set, so that any alignment between recovered
clusters and the independently computed rigor tiers is an empirical result rather than an
artifact of shared construction.

\paragraph{Feature set and preprocessing.}
Each school is represented by ten features spanning three axes: geography (one-hot indicators
for U.S.\ Census region), academic profile (the rigor classifier's opportunity-side component
scores---AP opportunity, CRDC advanced coursework, and test participation---plus the four-year
adjusted-cohort graduation rate), and resources (state--local funding per school-age resident
child and the county child-poverty rate). The two exam-performance components introduced in the
Week-5 rigor revision are excluded here: their coverage (approximately 31\%) would decimate the
complete-case population. Continuous features are $z$-standardized so that no single scale
dominates the Euclidean geometry. Consistent with the rest of the pipeline, we do not impute:
clustering is fit on the 5{,}801 schools (16.9\% of the 34{,}392-school universe) with complete
data on all ten features. This complete-case population is not a random subset---it requires a
school to appear in the federal funding, poverty, and academic collections simultaneously---and
all cluster-level findings are reported as statements about this subset, not the full universe.

\paragraph{Dimensionality and collinearity.}
Because the features are intercorrelated (funding correlates $0.68$ with the Northeast
indicator; poverty anti-correlates with funding and each academic component), we first apply
principal component analysis \citep{jolliffe2016}, both to quantify the collinearity that the
composite-indicator literature warns about \citep{nardo2008} and to obtain orthogonal axes for
visualization. Six components are required to reach 90\% cumulative variance (91.8\%),
indicating correlated but not redundant features. The leading components are interpretable:
PC1 (33.6\% of variance) is a \emph{regional wealth} axis loading positively on funding and
negatively on poverty; PC2 (17.6\%) is an \emph{academic intensity} axis loading on test
participation, advanced coursework, and AP opportunity. That the wealth axis explains nearly
twice the variance of the academic axis is itself a substantive description of how these
schools differ.

\paragraph{Cluster estimation and model selection.}
We fit $k$-means \citep{macqueen1967} for $k \in \{2,\dots,8\}$ and select $k$ with the gap
statistic under the one-standard-error rule \citep{tibshirani2001}, which selects $k=4$. As a
pre-registered cross-check we also compute average silhouette widths \citep{rousseeuw1987},
which peak at $k=2$ ($0.270$) and decline thereafter ($0.187$ at $k=4$). The two criteria
disagree, and we report the disagreement rather than resolving it in favor of the more
convenient answer; $k=4$ is retained because the gap statistic is the report's stated primary
criterion. As a robustness check on the partition itself, we re-cluster with Ward
agglomerative linkage \citep{ward1963} and compare partitions via the adjusted Rand index
\citep{hubert1985}: ARI $=0.39$, indicating broadly similar but non-identical structure.

\paragraph{Interpretation discipline.}
The best silhouette attained anywhere in the search ($0.270$) sits at the conventional
threshold for meaningful structure, and the retained $k=4$ solution falls below it. Following
the project's reporting standard, we present this as a finding about the data rather than a
modeling failure: on these features, U.S.\ high schools form a continuum with soft, overlapping
groupings, not well-separated types---which the moderate cross-algorithm agreement corroborates.
Downstream narratives therefore treat cluster labels as descriptive segments, not as ground
truth. Examined against the independently constructed rigor tiers, the clusters show real but
partial alignment: approximately 45\% of rigor-tier variance lies between clusters. The
patterning is informative in both directions. Two clusters reach high mean rigor through very
different resource profiles (mean funding of \$27{,}487 versus \$13{,}898 per resident child),
demonstrating that the recovered structure is not a simple wealth gradient; at the same time,
the lowest-rigor cluster is also the highest-poverty one (21.2\%), exactly the
socioeconomic-ordering pattern the composite-indicator literature cautions composite measures
tend to reproduce, and we flag it accordingly.

% ---------------------------------------------------------------------
% BibTeX for the NEW keys
% ---------------------------------------------------------------------
% @inproceedings{macqueen1967,
%   author    = {MacQueen, James},
%   title     = {Some Methods for Classification and Analysis of Multivariate Observations},
%   booktitle = {Proceedings of the Fifth Berkeley Symposium on Mathematical Statistics and Probability},
%   volume    = {1},
%   pages     = {281--297},
%   year      = {1967},
%   publisher = {University of California Press}
% }
% @article{ward1963,
%   author  = {Ward, Joe H.},
%   title   = {Hierarchical Grouping to Optimize an Objective Function},
%   journal = {Journal of the American Statistical Association},
%   volume  = {58},
%   number  = {301},
%   pages   = {236--244},
%   year    = {1963}
% }
% @article{tibshirani2001,
%   author  = {Tibshirani, Robert and Walther, Guenther and Hastie, Trevor},
%   title   = {Estimating the Number of Clusters in a Data Set via the Gap Statistic},
%   journal = {Journal of the Royal Statistical Society: Series B},
%   volume  = {63},
%   number  = {2},
%   pages   = {411--423},
%   year    = {2001}
% }
% @article{rousseeuw1987,
%   author  = {Rousseeuw, Peter J.},
%   title   = {Silhouettes: A Graphical Aid to the Interpretation and Validation of Cluster Analysis},
%   journal = {Journal of Computational and Applied Mathematics},
%   volume  = {20},
%   pages   = {53--65},
%   year    = {1987}
% }
% @article{hubert1985,
%   author  = {Hubert, Lawrence and Arabie, Phipps},
%   title   = {Comparing Partitions},
%   journal = {Journal of Classification},
%   volume  = {2},
%   number  = {1},
%   pages   = {193--218},
%   year    = {1985}
% }
% @article{jolliffe2016,
%   author  = {Jolliffe, Ian T. and Cadima, Jorge},
%   title   = {Principal Component Analysis: A Review and Recent Developments},
%   journal = {Philosophical Transactions of the Royal Society A},
%   volume  = {374},
%   number  = {2065},
%   pages   = {20150202},
%   year    = {2016}
% }
